\chapter{主理想整环上的模}\label{ch:b3:modules}

以\emph{环}代替域的线性代数：这一微小的
假设改动，便产生了代数学中最重要的统一
定理之一。$\Z$ 上的模即阿贝尔群；
$K[X]$ 上的模即配有自同态的向量空间。
主理想整环上有限生成模的结构定理
因此一举分类了全部有限生成
阿贝尔群\emph{与}相似意义下的全部自同态 ---
若尔当约化，本科二年级借助精巧的归纳法得到，
作为推论自然得出，连同它更精细的姊妹结果
\emph{有理}标准形，它在任意域上都成立。计算的核心工具是
史密斯标准形，
一种堪比欧几里得的矩阵算术。

全文中，$A$ 是交换环，之后将为主理想整环；“模”
指 $A$-模。

\section{模，自由模}

\begin{definition}\label{def:b3:modules:module}
称 \emph{$A$-模}\index{模} 为阿贝尔群 $(M, +)$，
它配有标量乘法 $A \times M \to M$，满足
向量空间公理：$a(x + y) = ax + ay$，$(a + b)x = ax + bx$，
$(ab)x = a(bx)$，$1x = x$。子模、商模 $M/N$、态射
（$A$-线性映射）、直和 $\bigoplus_i M_i$，以及
同构定理，其定义与证明与向量空间、阿贝尔群
的情形完全相同；特别地，有 $M/\ker f \cong
\operatorname{im} f$，其中 $f$ 为态射。
\end{definition}

\begin{example}\label{ex:b3:modules:examples}
三个典型例子。
\begin{enumerate}
  \item 设 $A = K$ 为域：模即向量空间。
  \item 设 $A = \Z$：模恰为阿贝尔群（$nx$ 必定为
        $x + \dots + x$），子模即子群。
  \item 设 $A = K[X]$：模即一个 $K$ 上的向量空间 $V$，连同
        一个 $K$-线性映射 $u\colon x \mapsto X\cdot x$ ---
        反之，对每个序对 $(V, u)$（其中 $u \in \mathcal L(V)$），
        便成为 $K[X]$-模，其中 $P \cdot x = P(u)(x)$。
        子模恰好是 $u$-稳定的子空间。
\end{enumerate}
$A$ 的理想恰为 $A$ 的子模；商环
$A/I$ 是 $A$-模。与向量空间不同，模可以具有
\emph{挠}：在 $\Z/6\Z$ 中，元素 $\bar 3 \ne 0$ 被
$2 \neq 0$ 零化。
\end{example}

\begin{definition}\label{def:b3:modules:free}
称 $M$ 为\emph{有限生成}\index{有限生成模}，
若 $M = Ax_1 + \dots + Ax_n$，其中 $x_i$ 为某些元素。
称 $M$ 为秩为 $n$ 的\emph{自由模}\index{自由模}，若 $M \cong A^n$，
即它有一组\emph{基}（一族生成元，
且它在 $A$ 上线性独立）。每个有限生成的 $M$ 都是
某个自由模的商：$(a_1, \dots, a_n) \mapsto \sum a_ix_i$
将 $A^n$ 满射到 $M$ 上。
\end{definition}

\begin{proposition}[秩的不变性]\label{prop:b3:modules:rank}
若 $A \neq 0$ 且 $A^m \cong A^n$，则 $m = n$。
\end{proposition}

\begin{proof}
取 $\mathfrak m$ 为 $A$ 的一个极大理想（见\cref{thm:b3:rings:krull}）
并令 $k = A/\mathfrak m$，它是一个域。设同构 $f \colon A^m
\to A^n$ 将 $\mathfrak m A^m$ 映入 $\mathfrak m A^n$
（由线性性），因此诱导商之间的同构
\[
A^m/\mathfrak m A^m \;\cong\; A^n/\mathfrak m A^n,
\qquad\text{即}\qquad k^m \cong k^n
\]
作为 $k$-向量空间（商 $A^m/\mathfrak m A^m$ 被
$\mathfrak m$ 零化，因此 $A$ 的作用经过 $k$；
标准基的像构成一组 $k$-基）。
域 $k$ 上的维数理论给出 $m = n$。
\end{proof}

\begin{theorem}[自由模的子模]\label{thm:b3:modules:submodule}
设 $A$ 为主理想整环，$M \subseteq A^n$ 为子模。则 $M$ 是
秩 $\leq n$ 的自由模。
\end{theorem}

\begin{proof}
对 $n$ 归纳。当 $n = 1$ 时：$M$ 是理想，故 $M = (0)$
（秩为 $0$ 的自由模），或 $M = dA \cong A$（$x \mapsto dx$ 是
单射，由整环的性质可得）。当 $n > 1$ 时，设 $\pi \colon A^n \to A$ 为
取最后一个坐标的映射。则 $\pi(M)$ 为理想，即 $(0)$ 或 $dA$。若
$(0)$，则 $M \subseteq A^{n-1} \times \{0\}$，对此归纳即可。
否则取 $x_0 \in M$ 使 $\pi(x_0) = d$。对每个 $x \in M$
唯一地写成
\[
x = \underbrace{\frac{\pi(x)}{d}}_{\in A}\, x_0 +
\Bigl(x - \tfrac{\pi(x)}d x_0\Bigr),
\qquad x - \tfrac{\pi(x)}d x_0 \in M \cap \ker\pi
\]
（$\pi(x) \in dA$，故系数属于 $A$）。于是 $M = Ax_0
\oplus (M \cap \ker \pi)$：和是直和，因为由 $\pi(ax_0) =
ad = 0$ 可推出 $a = 0$。由归纳法，$M \cap \ker\pi \subseteq
\ker \pi \cong A^{n-1}$ 为秩 $\leq n - 1$ 的自由模；添上
$x_0$（如上所见，它与 $\ker\pi$ 线性独立）后，得到
$M$ 的一组基，其基数 $\leq n$。
\end{proof}

\begin{remark}\label{rem:b3:modules:presentation}
由此可知，在主理想整环上，每个有限生成模 $M$ 都有一个
\emph{有限表示}：满射 $\varphi\colon A^n \to M$
其核是自由模，基为 $c_1, \dots, c_k$（$k \leq n$），且 $M
\cong A^n / \operatorname{im}(C)$，其中 $C \in M_{n,k}(A)$ 是
以 $c_j$ 为列的矩阵。理解 $M$ 就等于理解
\emph{$A$ 上的矩阵}在源与目标中更换基
之下的等价类 —— 这正是下一节的主题。
\end{remark}

\section{史密斯标准形}

\begin{definition}\label{def:b3:modules:equivalence}
称 $B, C \in M_{n,k}(A)$ 为\emph{等价}的，若 $C =
QBP$，其中 $Q \in GL_n(A)$，$P \in GL_k(A)$（
\emph{在 $A$ 上}可逆：行列式属于 $A^\times$）。等价的
表示矩阵定义同构的模 $A^n/
\operatorname{im}$（在 $A^n$ 与 $A^k$ 中更换基）。
\end{definition}

\begin{theorem}[史密斯标准形]\label{thm:b3:modules:smith}
设 $A$ 为主理想整环，$B \in M_{n,k}(A)$。则 $B$ 等价于
一个对角矩阵\index{史密斯标准形}
\[
\operatorname{diag}(d_1, d_2, \dots, d_r, 0, \dots, 0),
\qquad d_1 \mid d_2 \mid \cdots \mid d_r \neq 0 ,
\]
且各 $d_i$ 在相伴意义下唯一：$d_1 \cdots d_i$ 是
$i \times i$ 子式（取自 $B$）的一个最大公因数（特别地，该
最大公因数是等价下的一个不变量）。$d_i$ 称为
$B$ 的\emph{不变因子}\index{不变因子}。
\end{theorem}

\begin{proof}
\emph{存在性。}若 $B = 0$，则结论成立。否则考虑理想
$(b)$ 的集合，其中每个理想由与 $B$ 等价的矩阵之元素生成；
由于 $A$ 是诺特环，选取矩阵 $B'$，使其与 $B$ 等价，
并取 $d$ 为 $B'$ 的某个元素，使 $(d)$ 在该集合中\emph{极大}。
借助行与列的交换，将 $d$ 移到 $(1,1)$ 位置。

\emph{断言：$d$ 整除 $B'$ 的每个元素。}首先考虑第 1 列，若
$b_{i1}$ 不是 $d$ 的倍数，令 $e = \gcd(d, b_{i1}) = ud +
vb_{i1}$（贝祖等式），于是 $(e) \supsetneq (d)$。$2 \times 2$
矩阵技巧：对第 $1$ 行与第 $i$ 行作用如下
\[
\begin{pmatrix} u & v \\ -\dfrac{b_{i1}}e & \dfrac de
\end{pmatrix}
\in GL_2(A)
\qquad
\Bigl(\det = \tfrac{ud + v b_{i1}}e = 1\Bigr)
\]
得到一个等价矩阵，其元素 $e$ 位于 $(1,1)$ 位置：
这与 $(d)$ 的极大性矛盾。因此 $d$ 整除第 $1$ 列，
同理对第 $1$ 行亦然。用第 1 行
与第 1 列的倍数相减即可消去它们：$B'$ 等价于 $\begin{pmatrix} d &
0\\ 0 & B''\end{pmatrix}$。接着，$d$ 整除每个元素 $b$（取自
$B''$）：将 $b$ 所在的行加到第 1 行（初等变换；
新的第一行含 $d$ 与 $b$ 类元素），并重复
列消去论证：若有非倍数，则会再次使理想
$(d)$ 严格增大。现对阶数归纳：$B''$ 的全部元素
均被 $d$ 整除，并具有史密斯标准形 $\operatorname{diag}(d_2,
\dots)$，其元素仍被 $d$ 整除（任一
$QB''P$ 的元素都是 $A$ 上 $B''$ 元素的线性组合）；令 $d_1 =
d$。

\emph{唯一性。}记 $D_i(B)$ 为所有 $i \times i$
子式的一个最大公因数。行、列变换，更一般地，
乘以\emph{任意}矩阵，也不会使最大公因数变小：
$i \times i$ 子式（取自 $QB$）是 $A$ 上 $B$ 的子式的线性组合
（柯西--比内展开；或直接地：$QB$ 的每一行都是
$B$ 各行的组合，而子式关于行是多重线性的）。
因此 $D_i(QBP)$ 与 $D_i(B)$ 互相整除：$D_i$ 是
等价不变量。在对角形上，非零的 $i
\times i$ 子式恰为 $i$ 个 $d_j$ 的乘积，并且
由整除性 $d_1 \mid \dots \mid d_r$，$d_1 \cdots d_i$ 恰为
最大公因数。于是在相差单位的意义下 $d_1 \cdots d_i = D_i(B)$，且 $d_i =
D_i/D_{i-1}$ 唯一确定。
\end{proof}

\begin{method}\label{met:b3:modules:smith}
在欧几里得整环（$\Z$、$K[X]$）上，史密斯约化是一个
算法 --- 无需极大性论证：把欧几里得度最小的
元素移到 $(1,1)$ 位置；若它不能整除
同行或同列的某元素，则作欧几里得除法，留下
严格更小的余项于该处 --- 把它换入并重新开始
（终止性：度量递减）；当它整除其整行与整列时，
则消去它们；若它不能整除某内部元素，则把该行
加到第 $1$ 行并重新开始；对内部块递归。
实际操作整数矩阵时：计算 $D_1 = \gcd$（元素的最大公因数），
$D_2$、\dots\ 借助子式计算（阶数较小时），或直接运行算法。
\end{method}

\begin{example}[史密斯约化的完整实例]
\label{ex:b3:modules:smithworked}
对 $M = \begin{pmatrix} 1 & 2 & 3\\ 4 & 5 & 6\\ 7 & 8 & 9
\end{pmatrix}$ 在 $\Z$ 上作史密斯约化。左上角元素 $1$ 整除所有元素：
消去其所在的行与列（$L_2 \leftarrow L_2 - 4L_1$，$L_3
\leftarrow L_3 - 7L_1$，再作 $C_2 \leftarrow C_2 - 2C_1$，$C_3
\leftarrow C_3 - 3C_1$）：
\[
M \sim \begin{pmatrix} 1 & 0 & 0\\ 0 & -3 & -6\\ 0 & -6 & -12
\end{pmatrix} .
\]
在内部块中，角元 $-3$ 整除全部元素：$L_3
\leftarrow L_3 - 2L_2$ 与 $C_3 \leftarrow C_3 - 2C_2$ 将其消为
$\operatorname{diag}(-3, 0)$。调整符号（把一行乘以
$-1$，这是合法的运算）：
\[
M \sim \operatorname{diag}(1, 3, 0),
\qquad
\Z^3/M\Z^3 \cong \Z/3\Z \times \Z .
\]
用行列式因子作交叉验证：$D_1 =
\gcd(\text{元素}) = 1$；每个 $2\times2$ 子式（取自 $M$）都是
$3$ 的倍数（例如 $\det\bigl(\begin{smallmatrix}1 & 2\\ 4
& 5\end{smallmatrix}\bigr) = -3$），且其中有一个等于 $-3$：$D_2 =
3$；$D_3 = \det M = 0$。因此 $d_1 = 1$，$d_2 = 3$，$d_3 = 0$：
结果一致。两点启示：零不变因子记录了
秩的下降（余核多出一个自由的 $\Z$ 直和项），而
整除链 $1 \mid 3 \mid 0$ 便是史密斯判据
--- 违反该链的对角约化（例如
$\operatorname{diag}(2, 3)$，粗心者可能
由 $\bigl(\begin{smallmatrix}2 & 0\\ 0 &
3\end{smallmatrix}\bigr)$ 过早停止而得到，正确的史密斯
标准形为 $\operatorname{diag}(1, 6)$，因为这里 $D_1 = 1$！）该约化尚未
完成。
\end{example}

\begin{omfigure}
\begin{tikzpicture}[scale=0.72]
  \clip (-3.4,-1.3) rectangle (4.4,4.3);
  \foreach \x in {-3,...,4}
    \foreach \y in {-1,...,4}
      \fill[black!40] (\x,\y) circle (1.6pt);
  \fill[omDef!12, opacity=0.7] (0,0) -- (1,3) -- (3,3) -- (2,0)
    -- cycle;
  \foreach \a/\b in {0/0, 2/0, 4/0, -2/0, 1/3, 3/3, -1/3, -3/3,
                     2/6, -2/-6, 0/6}
    \fill[omThm] (\a,\b) circle (2.6pt);
  \draw[omThm, thick, ->] (0,0) -- (2,0)
    node[midway, below, font=\small] {$(2,0)$};
  \draw[omThm, thick, ->] (0,0) -- (1,3)
    node[midway, left, font=\small] {$(1,3)$};
\end{tikzpicture}

{\small $L = \Z(2,0) + \Z(1,3)$ 为 $\Z^2$ 的子格（红
点）。矩阵 $\bigl(\begin{smallmatrix} 2 &
1\\ 0 & 3\end{smallmatrix}\bigr)$ 的史密斯标准形为
$\operatorname{diag}(1, 6)$：在适配基 $f_1 = (1,3)$，
$f_2 = (0,1)$ 构成 $\Z^2$ 的基，则有 $L = \Z f_1 \oplus \Z\,6f_2$，
故 $\Z^2/L \cong \Z/6\Z$ --- 指数等于 $\abs{\det} = 6$，
即阴影基本区域的面积。}
\end{omfigure}

\section{结构定理}

\begin{definition}\label{def:b3:modules:torsion}
设 $A$ 为整环，$M$ 为 $A$-模。称下式所定义的
\emph{挠子模}\index{挠子模}为
\[
T(M) = \{x \in M : ax = 0 \text{ 对某个 } a \neq 0\}
\]
（子模：若 $ax = by = 0$，则 $ab(x + y) = 0$，$ab \ne 0$）。
$M$ 称为\emph{无挠}的，若 $T(M) = 0$，
若 $T(M) = M$，则称为\emph{挠模}。
\end{definition}

\begin{theorem}[主理想整环上有限生成模的
结构]\label{thm:b3:modules:structure}
设 $A$ 为主理想整环，$M$ 为有限生成的 $A$-模。存在
唯一的 $r \in \N$ 及非零非单位元 $d_1 \mid d_2 \mid
\cdots \mid d_s$，在相伴意义下唯一，且
使得\index{模的结构定理}
\[
M \;\cong\; A^r \,\oplus\, A/(d_1) \oplus \cdots \oplus A/(d_s).
\]
此外 $T(M) \cong A/(d_1)\oplus\dots\oplus A/(d_s)$，且 $M/T(M)
\cong A^r$：主理想整环上有限生成的\emph{无挠}模
是自由模。
\end{theorem}

\begin{proof}
\emph{存在性。}写出 $M \cong A^n/\operatorname{im}(C)$
（见\cref{rem:b3:modules:presentation}），并将 $C$ 化为史密斯标准形：
经两次基变换后，$M \cong A^n / (d_1A \times \dots
\times d_rA \times 0 \times \dots \times 0) \cong A/(d_1) \oplus
\dots \oplus A/(d_r) \oplus A^{\,n - r}$。舍去
$d_i$ 为单位的因子（$A/(d_i) = 0$）；
整除链仍然成立。

\emph{挠的识别。}在该分解中，$A^r$ 是
无挠的（整环无零因子），且每个 $A/(d_i)$
是挠的（被 $d_i \ne 0$ 零化）；直和相应地分解挠，
即 $T(M) = \bigoplus_i A/(d_i)$，且 $M/T(M) \cong A^r$。

\emph{$r$ 的唯一性}：$M/T(M) \cong A^r$ 只依赖于 $M$，
而由\cref{prop:b3:modules:rank}确定 $r$。

\emph{$d_i$ 的唯一性}：只需讨论挠模
$T = T(M)$。将每个 $d_i$ 分解为素元之积，并
用中国剩余定理（\cref{thm:b3:rings:crt}；不同素元
生成互素理想）加以拆分：
\[
A/(d) \cong \bigoplus_{p \mid d} A/\bigl(p^{v_p(d)}\bigr):
\qquad
T \cong \bigoplus_{p} \bigoplus_{j} A/\bigl(p^{k_{p,j}}\bigr),
\]
即\emph{初等因子}\index{初等因子}
$p^{k_{p,j}}$。反之，$d_i$ 可由初等因子的
重集重构（$d_s$ 为各素元最高幂之积，依此类推），
故只需证明
重集 $\{k_{p,j}\}_j$ 由 $T$ 决定，对每个素元
$p$。固定 $p$；对 $j \geq 1$，考虑 $A/(p)$-向量空间
$p^{j-1}T/p^jT$。对循环因子 $A/(p^k)$：
\[
p^{j-1}\bigl(A/(p^k)\bigr)\big/p^{j}\bigl(A/(p^k)\bigr)
\cong
\begin{cases}
A/(p) & \text{若 } j \leq k,\\
0 & \text{若 } j > k,
\end{cases}
\]
而在因子 $A/(q^k)$（$q \neq p$）上，乘以 $p$ 为
双射（$p$ 模 $q^k$ 可逆，由贝祖定理），故商
为 $0$。直和逐项成立：$\dim_{A/(p)}
p^{j-1}T/p^jT = \#\{i : k_{p,i} \geq j\}$。这些内在的
维数决定指数的重集。
\end{proof}

\begin{corollary}[有限生成交换群]
\label{cor:b3:modules:abelian}
每个有限生成交换群同构于 $\Z^r \times
\Z/d_1\Z\times\dots\times\Z/d_s\Z$，其中 $d_1 \mid \dots \mid
d_s$，唯一确定。每个\emph{有限}交换群都是若干
素数幂阶循环群之积，且作为
重集唯一。\index{交换群，结构}
\end{corollary}

\begin{example}\label{ex:b3:modules:abelian}
阶为 $p^n$ 的交换群与 $n$ 的\emph{分拆}一一对应。
对 $p^4$，有：$\Z/p^4$，$\Z/p^3\times\Z/p$，
$(\Z/p^2)^2$，$\Z/p^2 \times (\Z/p)^2$，$(\Z/p)^4$ --- 共五个群，
因为 $4$ 有五种分拆。混合阶的情形可逐个素数相乘计数
（由中国剩余定理），共有 $5 \times 2$ 个交换群，
其阶为 $2^4 \cdot 3^2 = 144$。
\end{example}

\section{应用：自同态的标准形}

设 $K$ 为域，$V$ 为有限维 $K$-向量空间，
维数为 $n$，$u \in \mathcal L(V)$；令 $V$ 成为 $K[X]$-模，其数乘为 $P
\cdot x = P(u)(x)$（见\cref{ex:b3:modules:examples}）。该模
是有限生成的（$K$-基生成它），并且是挠模：
对每个 $x$，$n+1$ 个向量 $x, u(x), \dots, u^n(x)$ 在 $K$ 上线性相关，
从而给出一个非零的零化多项式。

\begin{definition}\label{def:b3:modules:companion}
设 $P = X^m + a_{m-1}X^{m-1} + \dots + a_0$ 为首一多项式，
称其\emph{友矩阵}\index{友矩阵}为：
\[
C_P = \begin{pmatrix}
0 & & & -a_0\\
1 & \ddots & & -a_1\\
& \ddots & 0 & \vdots\\
& & 1 & -a_{m-1}
\end{pmatrix} :
\]
它是“乘以 $X$”这一运算在 $K[X]/(P)$ 上，
关于基 $1, \bar X, \dots, \bar X^{m-1}$ 的矩阵。
\end{definition}

\begin{theorem}[弗罗贝尼乌斯：有理标准形]
\label{thm:b3:modules:frobenius}
存在唯一的首一非常数多项式序列 $P_1
\mid P_2 \mid \dots \mid P_s$（
\emph{相似不变量}\index{相似不变量}，由 $u$ 确定），使得
在 $K[X]$-模的意义下，
\[
V \cong K[X]/(P_1) \oplus \cdots \oplus K[X]/(P_s):
\]
在适当的基下，$u$ 的矩阵为分块对角矩阵
$\operatorname{diag}(C_{P_1}, \dots, C_{P_s})$。此外：
\begin{enumerate}
  \item $P_s = \mu_u$（极小多项式），且 $P_1\cdots P_s =
        \chi_u$（特征多项式）；特别地，$\mu_u
        \mid \chi_u$ \emph{（凯莱–哈密顿定理的重新证明）}，且
        $\chi_u \mid \mu_u^{\,s}$，因此 $\chi_u$ 与 $\mu_u$ 有
        相同的不可约因子。
  \item 两个自同态（或方阵）相似，当且仅当
        它们有相同的相似不变量。
\end{enumerate}
\end{theorem}

\begin{proof}
将结构定理（\cref{thm:b3:modules:structure}）
应用于主理想整环 $K[X]$：挠模 $V$ 分解，
不变因子为 $P_i$，并规范化为首一多项式（$K[X]$ 的单位为 $K^\times$）；
不出现自由部分（$V$ 为挠模）。在每个循环
因子 $K[X]/(P_i)$ 上，乘以 $X$ 的矩阵为 $C_{P_i}$，
在 $\bar X$ 的幂基下；将各基拼接，
即得分块形式。

(1) $V = \bigoplus K[X]/(P_i)$ 的零化子为 $(P_1)\cap
\dots\cap(P_s) = (P_s)$（整除链：$P_s$ 为公共倍式，
且最后一个因子中 $1$ 的类恰被 $(P_s)$ 零化）：$\mu_u = P_s$。
对于 $\chi_u$：在循环因子上，$\chi_{C_P} = P$，
对 $m = \deg P$ 作归纳。沿第一行展开
$\det(XI_m - C_P)$（第一行的元素依次为 $X$、
零，最后一列为 $a_0$）：
\[
\det(XI_m - C_P) = X\,\det\bigl(XI_{m-1} - C_{\tilde P}\bigr)
+ (-1)^{1+m}\,a_0\,\det L ,
\]
其中 $\tilde P = X^{m-1} + a_{m-1}X^{m-2} + \dots + a_1$（形状相同，
次数低一），$L$ 为对角线元素为 $(-1, \dots, -1)$ 的三角矩阵，
故 $\det L = (-1)^{m-1}$。由归纳假设，第一项为 $X\tilde P$，
第二项为 $a_0$，故总和为
$X\tilde P + a_0 = P$（基础情形 $m=1$：$\det(X + a_0) = P$）。
行列式在分块上相乘：
$\chi_u = \prod P_i$。凯莱–哈密顿：$\chi_u \in (\mu_u)$，
由于 $P_s \mid$ 每个\dots\ 反之每个 $P_i \mid P_s$，故
$\chi_u = \prod P_i$ 整除 $P_s^{\,s} = \mu_u^s$；而 $\mu_u =
P_s$ 作为其因子整除 $\chi_u$。

(2) 相似的自同态即为共轭的模结构，因此
有相同的不变量（由
\cref{thm:b3:modules:structure} 的唯一性）；反之，相同的不变量
给出同构的 $K[X]$-模，而模同构恰为
使两个自同态相互协调的线性双射，即
相似。
\end{proof}

\begin{corollary}[相似性对域扩张不敏感]
\label{cor:b3:modules:descent}
设 $K \subseteq L$ 为域，$M, N \in M_n(K)$。若 $M$ 与
$N$ 在 $L$ 上相似，则它们在 $K$ 上也相似。
\end{corollary}

\begin{proof}
$M$ 的相似不变量可由史密斯子式公式
（\cref{thm:b3:modules:smith}）计算，该公式应用于表示矩阵
$XI_n - M$（在 $K[X]$ 上）——事实上，$K[X]$-模
$V_M = K^n$ 的表示矩阵为 $XI_n - M$：映射 $K[X]^n \to V_M$，
$(Q_i) \mapsto \sum Q_i(M)e_i$ 是满射，核由
$XI_n - M$ 的各列生成（直接验证：模去这些列后，$K[X]^n$ 中的每个元素都化为常值向量，
而常值向量双射地映到该模上）；周末问题对此作了详细说明。
多项式的最大公因数在域扩张下不变：
若 $d$ 为 $K[X]$ 中族 $(f_j)$ 的首一最大公因数，
由贝祖等式有 $d = \sum u_jf_j$，$u_j \in K[X]$，
故 $f_j$ 在 \emph{$L[X]$ 中}的任一公因数都整除 $d$；
而 $d$ 本身也是公因数，
因此它也是 $L[X]$ 中的最大公因数。
于是，$XI_n - M$ 的不变因子（即相继子式最大公因数之商）在
$K$ 与 $L$ 上相同：
$M, N$ 在 $L$ 上有相同的相似不变量当且仅当它们在 $K$ 上有相同的相似不变量；
由\cref{thm:b3:modules:frobenius}(2) 即得结论。
\end{proof}

\begin{theorem}[若尔当标准形（重新推导）]\label{thm:b3:modules:jordan}
设 $\chi_u$ 在 $K$ 上分裂（例如 $K = \C$）。对
$V$ 作初等因子分解（
见\cref{thm:b3:modules:structure}的证明），而不用不变因子：
\[
V \cong \bigoplus_{\lambda, j} K[X]\big/\bigl((X -
\lambda)^{k_{\lambda,j}}\bigr),
\]
在每个因子的基 $\bigl(\overline{(X-\lambda)^{k-1}}, \dots,
\overline{(X - \lambda)}, \bar 1\bigr)$ 下，$u$ 作用为
若尔当块 $J_k(\lambda)$：每个特征多项式分裂的自同态
都有若尔当基，且块的多重集
$(\lambda, k)$ 是唯一的。\index{若尔当标准形}
\end{theorem}

\begin{proof}
挠模 $V$ 的初等因子是
$(X - \lambda)^k$，其中 $X - \lambda$ 取遍
$\mu_u$ 的不可约因子（它也分裂，因为 $\chi_u$ 分裂，且两者
有相同的不可约因子，见
\cref{thm:b3:modules:frobenius}）。在 $W = K[X]/((X-\lambda)^k)$ 中，
令 $f_j = \overline{(X - \lambda)^{k-j}}$，$j = 1, \dots, k$：
则 $(X - \lambda)f_j = f_{j-1}$（$f_0 = 0$），即 $u(f_j)
= \lambda f_j + f_{j-1}$：$u$ 在 $(f_1, \dots,
f_k)$ 上的矩阵恰为 $J_k(\lambda)$（对角线上方为 1）。
初等因子多重集的唯一性见
\cref{thm:b3:modules:structure}。
\end{proof}

\begin{remark}\label{rem:b3:modules:overview}
标准形的层次结构现在一目了然：
\emph{有理}标准形在每个域上都存在，并且
绝对地刻画相似性（见\cref{cor:b3:modules:descent}）；
\emph{若尔当} 标准形是其在 $\chi_u$ 分裂时的精细化。
本科二年级用维数计数给出的若尔当定理证明
被包含在内：所有组合计数不过是主理想整环
$K[X]$ 上的算术。
\end{remark}

\section{练习}

\begin{exercise}[$\star$]\label{exo:b3:modules:1}
(a) 证明 $\Q$ 作为 $\Z$-模不是有限生成的。
(b) 证明 $\Q$ 无挠但不是自由的。
(c) 为什么这两个结论都不与
\cref{thm:b3:modules:structure}矛盾？
\end{exercise}

\begin{exercise}[$\star$]\label{exo:b3:modules:2}
在同构意义下列出阶为 $360$ 的所有交换群，
分别写出初等因子形式和不变因子形式。
阶为 $p^5$ 的交换群有多少个？
\end{exercise}

\begin{exercise}[$\star$]\label{exo:b3:modules:3}
求下列矩阵在 $\Z$ 上的史密斯标准形：
\[
B = \begin{pmatrix} 2 & 4\\ 6 & 8 \end{pmatrix},
\qquad
C = \begin{pmatrix} 2 & 0 & 0\\ 0 & 3 & 0\\ 0 & 0 & 12
\end{pmatrix},
\]
并确定交换群 $\Z^2/B\Z^2$ 与 $\Z^3/C\Z^3$。
\end{exercise}

\begin{exercise}[$\star\star$]\label{exo:b3:modules:4}
设 $L \subseteq \Z^n$ 是秩为 $n$ 的子群，其基由
$B \in M_n(\Z)$ 的列构成，$\det B \neq 0$。证明
$\Z^n/L$ 是有限的，基数为 $\abs{\det B}$，并证明 $\Z^n/L
\cong \prod_i \Z/d_i\Z$，其中 $d_i$ 为 $B$ 的不变因子。
以 $L = \Z(2,0) + \Z(1,3)$ 为例说明。
\end{exercise}

\begin{exercise}[$\star\star$]\label{exo:b3:modules:5}
设 $A$ 为整环。(a) 验证 $T(M)$ 是子模，且
$M/T(M)$ 无挠。(b) 设 $(X, Y)$ 为 $K[X,Y]$ 的理想，证明它作为
$K[X,Y]$-模是无挠的，但不是自由的：
结构定理确实需要主理想整环这一假设。
\end{exercise}

\begin{exercise}[$\star\star$]\label{exo:b3:modules:6}
(a) 证明 $2\Z$ 在 $\Z$-模
$\Z$ 中没有直和补：自由模的子模是自由的
（见\cref{thm:b3:modules:submodule}），但直和因子
不必是自由的。(b) 证明：若 $M \subseteq A^n$（$A$ 为主理想整环）满足：
$A^n/M$ 无挠，则 $M$ \emph{是}直和因子。
\end{exercise}

\begin{exercise}[$\star\star$]\label{exo:b3:modules:7}
在 $\Z^2$ 中求解方程组
\[
\begin{cases}
2x + 4y \equiv b_1 \pmod{20},\\
6x + 8y \equiv b_2 \pmod{20},
\end{cases}
\]
并确定对哪些数对 $(b_1, b_2)$ 方程组有解，利用
\cref{exo:b3:modules:3}中的史密斯标准形（对
两侧作可逆变量替换）。
\end{exercise}

\begin{exercise}[$\star\star$]\label{exo:b3:modules:8}
(a) 确定幂零 $4 \times 4$ 矩阵的所有相似不变量及可能的
若尔当标准形，按它们所实现的
$4$ 的分拆排序。
(b) 给出两个 $4\times4$ 复矩阵，它们具有相同的
特征多项式\emph{和}极小多项式，但不相似，
并证明当 $n \leq 3$ 时这种情形不可能出现。
\end{exercise}

\begin{exercise}[$\star\star\star$]\label{exo:b3:modules:9}
设 $u \in \mathcal L(V)$，$\dim V = n$。证明下列各条
陈述等价：(i) $V$ 是循环 $K[X]$-模（存在 $x$
使得 $V = K[u]x$，称\emph{循环向量}）；(ii) $\mu_u = \chi_u$；
(iii) 在\cref{thm:b3:modules:frobenius}中 $s = 1$。由此推出
友矩阵有循环向量，并确定何时
对角矩阵有循环向量。
\end{exercise}

\begin{exercise}[$\star\star\star$]\label{exo:b3:modules:10}
设 $M \in M_n(\Z)$，视为 $\Z^n$ 上的自同态，证明指标公式：
若 $\det M \ne 0$，则
$[\Z^n : M\Z^n] = \abs{\det M}$，并由此推出 $M \in GL_n(\Z)$
当且仅当 $\det M = \pm 1$。应用：群 $\Z^2$ 恰有
$\sigma_1(m) = \sum_{d \mid m} d$ 个指数为 $m$ 的子群。
\emph{（计数埃尔米特标准形的矩阵 $\bigl(\begin{smallmatrix}
a & b\\ 0 & d\end{smallmatrix}\bigr)$，$ad = m$，$0 \leq b <
d$。）}
\end{exercise}

\begin{exercise}[$\star\star$]\label{exo:b3:modules:11}
（阿贝尔群中的方程 $x^k = e$）设 $G$ 为有限
阿贝尔群，其不变因子为 $d_1 \mid d_2 \mid \dots
\mid d_s$。
(a) 证明对每个 $k \geq 1$，
\[
\#\{x \in G : x^k = e\} \;=\; \prod_{i=1}^{s}\gcd(k, d_i) .
\]
(b) 由此推出：有限阿贝尔群是循环群，当且仅当
对每个 $k$，方程 $x^k = e$ 至多有 $k$
个解。
(c) 重新推出 $K^\times$ 的有限子群的循环性
（$K$ 为域，见\cref{ch:b3:galois}）：为何多项式
$X^k - 1$ 能保证 (b) 中的判别准则成立？
\end{exercise}

\begin{exercise}[$\star\star\star$]\label{exo:b3:modules:12}
（初等矩阵生成）(a) 证明 $M \in M_n(\Z)$
在 $M_n(\Z)$ 中可逆，当且仅当 $\det M = \pm1$。
(b) 证明 $SL_2(\Z)$ 由两个初等
矩阵生成，它们分别是
$E = \bigl(\begin{smallmatrix}1 & 1\\ 0 &
1\end{smallmatrix}\bigr)$ 与
$F = \bigl(\begin{smallmatrix}1 & 0\\ 1 &
1\end{smallmatrix}\bigr)$。
\emph{（对 $M
\in SL_2(\Z)$ 的第一列作欧几里得算法，左乘 $E, F$ 的幂，
化为 $\pm\bigl(\begin{smallmatrix}1 & *\\ 0 &
1\end{smallmatrix}\bigr)$；手算收尾 --- 注意
$-I = (EF^{-1}E)^2$。）}
(c) 说明与史密斯标准形约化的联系：在 $\Z$ 上，
行列式为 $1$ 的行与列变换足以
化为对角形，至多相差符号。
\end{exercise}

\section{问题：交换化子与双交换化子}

\begin{problem}[{周末问题 --- 有理形式、交换化子，
双交换化子}]\label{pb:b3:modules:1}
设 $K$ 为域，$V$ 是 $K$ 上维数为 $n \geq
1$ 的向量空间，$u \in \mathcal L(V)$。我们研究\emph{交换化子}
\[
\mathcal C(u) = \{v \in \mathcal L(V) : uv = vu\},
\]
它是 $\mathcal L(V)$ 的子代数，且包含 $K[u] = \{P(u) : P \in
K[X]\}$，并证明弗罗贝尼乌斯维数公式与双
交换化子定理：$\mathcal C(\mathcal C(u)) = K[u]$。全文中，
$V$ 是 $K[X]$-模（由 $u$ 定义），其不变因子为
$P_1 \mid \cdots \mid P_s$，循环分解为 $V =
\bigoplus_{i=1}^s V_i$，$V_i = K[u]\,x_i \cong K[X]/(P_i)$，$n_i
= \deg P_i$（\cref{thm:b3:modules:frobenius}）。

\medskip
\noindent\textbf{第一部分 --- 表出矩阵 $XI - M$，与
热身题。}
\begin{enumerate}
  \item 设 $M \in M_n(K)$，令 $\varphi \colon K[X]^n \to V_M
        = K^n$ 把 $(Q_1, \dots, Q_n)$ 映为 $\sum_i Q_i(M)e_i$。
        证明 $\varphi$ 为满的
        $K[X]$-模同态，且 $XI_n - M$ 的每一列
        都在 $\ker\varphi$ 中。
  \item 证明：模 $XI_n - M$ 的各列，
        $K[X]^n$ 中每个元素都同余于一个常向量
        \emph{（利用 $Xe_i \equiv \sum_j m_{ji}e_j$ 降低次数）}，
        由此推出 $\ker\varphi = (XI_n - M)\,K[X]^n$，模
        $V_M$ 的表出矩阵为 $XI_n - M$。恢复
        \cref{cor:b3:modules:descent} 的出发点：
        $M$ 的相似不变量是
        $XI_n - M$ 的非单位不变因子。
  \item 计算下列矩阵的相似不变量：数量矩阵
        $\lambda I_n$；对角元两两互异的
        对角矩阵；$n \times n$ 的若尔当块 $J_n(0)$；
        $\operatorname{diag}(J_2(0), J_1(0))$，取 $n = 3$。
  \item 证明 $\dim K[u] = \deg \mu_u = n_s$。
\end{enumerate}

\medskip
\noindent\textbf{第二部分 --- 循环模之间的同态。}
\begin{enumerate}[resume]
  \item 设 $P, Q$ 为首一的非常数多项式。证明
        $K[X]$-同态 $f \colon K[X]/(P) \to K[X]/(Q)$
        由 $f(\bar 1)$ 唯一确定，且 $\bar c \in K[X]/(Q)$
        可作为 $f(\bar 1)$，当且仅当 $P\bar c = 0$ 在 $K[X]/(Q)$ 中。
  \item 由此推出
        \[
        \operatorname{Hom}_{K[X]}\bigl(K[X]/(P),\, K[X]/(Q)\bigr)
        \;\cong\; K[X]\big/\bigl(\gcd(P, Q)\bigr),
        \]
        其维数为 $\deg \gcd(P, Q)$（以 $K$ 为基域）。\emph{（证明
        $\bar c$ 为满足 $P\bar c = 0$（在 $K[X]/(Q)$ 中）的解
        构成由
        $\overline{Q/\gcd(P,Q)}$ 生成的循环子模。）}
  \item 证明\emph{弗罗贝尼乌斯公式}：
        \[
        \dim_K \mathcal C(u) = \sum_{i,j=1}^{s} \deg\gcd(P_i,
        P_j) = \sum_{i=1}^{s} (2s - 2i + 1)\, n_i .
        \]
        \emph{（可交换的 $v$ 恰为 $K[X]$ 在
        $V$ 上的自同态；将 $\operatorname{End}(\bigoplus_i V_i)$ 分解为
        以 $V_j \to V_i$ 的同态为元素的矩阵，并利用
        整除链。）}
  \item 由此推出 $\dim \mathcal C(u) \geq n$，且等号成立当且仅当 $u$
        为循环的（$s = 1$），并计算 $\dim\mathcal C(u)$
        在 $u = \lambda\,\mathrm{id}$ 时的值，即公式的两个极端情形。
  \item 直接验证弗罗贝尼乌斯公式在
        $\operatorname{diag}(J_2(0), J_1(0))$ 上成立，方法是
        显式计算其交换化子（写成 $3\times3$ 矩阵）。
\end{enumerate}

\medskip
\noindent\textbf{Part III --- 双重中心化子定理。}
设 $w \in \mathcal C(\mathcal C(u))$，证明 $w \in K[u]$。
\begin{enumerate}[resume]
  \item 证明 $K[u] \subseteq \mathcal C(\mathcal C(u))$，并且
        每个 $w \in \mathcal C(\mathcal C(u))$ 都与 $u$ 可交换
        --- 因此待证的包含关系
        $\mathcal C(\mathcal C(u)) \subseteq K[u]$ 是对
        $w \in \mathcal C(u)$ 的真正加强。
  \item 首先设 $u$ 是\emph{循环}的，即 $V = K[u]x$。
        直接证明 $\mathcal C(u) = K[u]$
        \emph{（将可交换的 $v$ 作用于 $x$：$v(x) = P(u)x$
        对某个 $P$ 成立，并比较 $v$ 与 $P(u)$
        在基 $u^k x$ 上的作用）}，从而在此情形下得证定理。
  \item 回到一般情形。对每个 $i$，令 $\pi_i\colon
        V \to V_i$ 为沿其余直和项的投影。
        证明 $\pi_i \in \mathcal C(u)$，并由此推出 $w$
        保持每个 $V_i$ 不变，且与 $u_i =
        u\restriction_{V_i}$ 可交换，于是
        对循环的 $u_i$ 应用第 11 题，存在多项式 $Q_i$ 使
        $w\restriction_{V_i} = Q_i(u)\restriction_{V_i}$。
  \item 剩下的是把各 $Q_i$ 粘合成一个多项式。对 $i
        \leq j$（从而 $P_i \mid P_j$），证明 $\eta_{ij} \colon
        V_j \to V_i$，$R(u)x_j \mapsto R(u)x_i$ 是
        良定义的 $K[X]$-同态 \emph{（需要检验的是
        由 $R(u)x_j = 0$ 可推出 $R(u)x_i = 0$）}，并且
        $\tilde\eta_{ij} = \eta_{ij}\circ\pi_j$ 以 $0$
        延拓到其余直和项上，且属于 $\mathcal C(u)$。
  \item 利用 $w\tilde\eta_{ij} = \tilde\eta_{ij}w$，证明 $Q_i
        \equiv Q_j \pmod{P_i}$（对 $i \leq j$）。由此推出 $Q =
        Q_s$
        满足 $Q \equiv Q_i \pmod {P_i}$（对所有 $i$），因此
        $w = Q(u)$ 在每个 $V_i$ 上成立，从而在 $V$ 上亦然：
        \[
        \boxed{\ \mathcal C(\mathcal C(u)) = K[u].\ }
        \]
  \item （尾声）由定理推出：若 $v$ 与
        \emph{每一个}与 $u$ 可交换的矩阵都可交换，且 $u$ 为
        循环的，则 $v$ 为 $u$ 的多项式；并给出一个
        例子说明 $\mathcal C(u) = K[u]$ 对 $u =
        \mathrm{id}$，$n \geq 2$ 时不成立 --- 循环性究竟
        在何处起作用？
\end{enumerate}

\medskip
\noindent\textbf{Part IV --- 相似不变量的
收获。}有理标准形是一台机器；以下列出它的
五个经典结论。
\begin{enumerate}[resume]
  \item （转置）证明每个 $M \in M_n(K)$ 都相似于
        其转置 ${}^tM$。\emph{（将 $XI - M$ 化为史密斯标准形的各步运算，
        经转置后，可将 $XI - {}^tM$ 化为
        同一个史密斯标准形；因此二者有相同的相似不变量。）}
  \item （相似性的下降）设 $K \subseteq L$ 为域
        扩张，$M, N \in M_n(K)$。证明若 $M$ 与 $N$
        在 $L$ 上相似，则它们在 $K$ 上亦相似。
        \emph{（$XI - M$ 在 $K[X]$ 中算得的史密斯标准形
        在 $L[X]$ 中仍为史密斯标准形 --- 为何
        不变因子不会改变？）}值得记住的推论：两个
        实矩阵在 $GL_n(\C)$ 中共轭，则它们在
        $GL_n(\R)$ 中亦共轭。
  \item （幂零分类）设 $u$ 为幂零算子。证明
        其分解为幂零若尔当块时，大小 $\geq k$ 的块个数
        等于
        $\operatorname{rk}u^{k-1} - \operatorname{rk}u^k$，并
        由此推出：$M_n(K)$ 中的幂零类（$K$ 为任意域）
        与 $n$ 的划分一一对应。
        $M_5(K)$ 中共有多少个幂零类？
  \item （一个具体例子）求
        导子 $D\colon P \mapsto P'$ 作用于
        $K_{n-1}[X]$（次数 $< n$ 的多项式空间）上的相似不变量：(a) 当 $K =
        \Q$ 时；(b) 当 $K = \mathbb F_p$ 且 $p < n$ 时
        \emph{（特征为 $p$ 时，$(X^p)' = 0$：计算
        $\ker D^k$，并利用第 18 题）}。
  \item （$GL_2(\mathbb F_q)$ 的共轭类）利用
        不变因子，证明
        $GL_2(\mathbb F_q)$ 的每个共轭类恰属于以下四种类型之一：
        中心型 $aI$；可对角化，具有两个不同的
        特征值 $a \neq b$，均属于 $\mathbb F_q^\times$；
        非半单，极小多项式为 $(X - a)^2$；
        循环型，特征多项式不可约。
  \item 计算每种类型的共轭类个数，并由此得出：
        $GL_2(\mathbb F_q)$ 恰有 $q^2 - 1$ 个
        共轭类。\emph{（计数首一不可约二次多项式，
        其系数属于 $\mathbb F_q$；无序对 $\{a, b\}$ 的个数；注意
        可逆性限制了常数项。）}
  \item （循环性是普遍的）证明 $M \in M_2(\mathbb F_q)$
        不是循环的，当且仅当 $M$ 为数量矩阵，并由此推出
        均匀随机选取的 $2\times2$ 矩阵（元素属于 $\mathbb F_q$）
        以概率 $1 - q^{-3}$ 是循环的。请陈述类似的
        对 $M_n$ 及大 $q$ 的启发式结论（无需证明）：
        非循环矩阵很少见 --- 这也正是
        \cref{pb:b3:modules:1} 的 Part III 只有在
        越过一般情形后才需要真正的工作。
\end{enumerate}

\medskip
\noindent\textbf{Part V --- 补充内容。}
\begin{enumerate}[resume]
  \item （中心化子的中心）证明
        代数 $\mathcal C(u)$ 的中心恰为 $K[u]$。\emph{（结合
        Part III 中的两个包含关系）}。由此推出
        $\mathcal C(u)$ 可交换，当且仅当 $u$ 是循环的 ---
        以纯结构的途径重新得到第 8 题中
        取等的情形。
  \item （哪些维数会出现？）由弗罗贝尼乌斯公式推出
        $\dim\mathcal C(u) \equiv n \pmod 2$ 对
        每个 $u$ 成立。然后确定 $\dim \mathcal C(u)$
        所取值的精确集合，当 $u$ 遍历
        $\mathcal L(V)$（$\dim V = 4$）时，并证明该集合为 $\{4, 6,
        8, 10, 16\}$。\emph{（列举满足 $n_1
        \leq \dots \leq n_s$ 且和为 $4$ 的度数序列，并用幂零算子实现每一种
        情形）}。特别地，$12$ 与 $14$ 虽奇偶性正确，
        却都取不到：奇偶性条件
        是必要的，但不是充分的。
  \item （$GL_2(\mathbb F_3)$ 的类方程）取 $q = 3$，
        计算第 20 题中每个共轭类的大小：
        利用轨道–稳定子定理：一个
        \emph{循环}矩阵 $M$ 在 $GL_2(\mathbb F_q)$ 中的中心化子
        是 $K[M]$ 的单位群（第 11 题）。确定
        $K[M]$ 在三种非中心类型中的形式，列出三个首一不可约
        二次多项式，其系数属于 $\mathbb F_3$，并验证
        类方程
        \[
        48 = \abs{GL_2(\mathbb F_3)}
        = 2\cdot1 + 1\cdot12 + 2\cdot8 + 3\cdot6,
        \]
        共 $2 + 1 + 2 + 3 = 8 = q^2 - 1$ 个类，
        与第 21 题的预言一致。
\end{enumerate}
\end{problem}
